\section{Introduction}

Quantum software is increasingly developed with familiar software-engineering workflows, but it is tested under an execution model that makes ordinary test oracles fragile. A quantum program usually returns samples from a distribution rather than a single deterministic value. The full quantum state is not directly observable without changing the computation, and the number of possible states grows exponentially with the number of qubits. Even when the intended circuit is correct, finite shots and backend noise can move observed frequencies enough to make exact-output assertions misleading. These properties turn regression testing into an oracle problem: developers need fault indicators without writing a fully specified expected distribution for every program and backend configuration~\cite{preskill2018quantum,huang2019statistical,li2020projection}.

A point that is easy to miss in this setting is \emph{where} the bug actually lives. A Qiskit program is not just the final \texttt{QuantumCircuit} object that reaches the backend. It is a builder that allocates qubits and classical registers, a set of parameterized APIs, a mapping from logical roles to wire indices, a measurement layout, and ancilla protocols that must uncompute cleanly. Many real defects are introduced at this program-contract layer---a permuted control qubit in a multi-controlled gate, a classical register read back in the wrong endianness, a QFT applied in the wrong direction, a rotation angle that breaks an intended periodicity, an ancilla that is never uncomputed---\emph{before} the concrete circuit is emitted. A metamorphic transformation that rewrites the \emph{final} circuit takes the buggy builder output as ground truth and rewrites around it; it therefore cannot expose a bug that is already baked into that output.

Existing work answers different parts of the oracle problem. Bugs4Q provides reproducible quantum-program defects for controlled studies~\cite{miao2023bugs4q}. MorphQ uses metamorphic transformations over generated circuits to expose faults in the Qiskit platform~\cite{paltenghi2023morphq}. Differential testing compares quantum software stacks when multiple implementations can run comparable programs~\cite{wang2021qdiff}. Assertion-based and projective-measurement approaches let developers encode stronger probabilistic or quantum-specific properties when those properties are known~\cite{huang2019statistical,li2020projection,oldfield2025faster}. These techniques operate at the \emph{circuit} or \emph{platform} level. We focus on the complementary boundary they leave open: the program/builder contracts that determine the circuit in the first place.

We target that boundary with QMTester, a metamorphic-testing workflow whose relations are defined over \emph{program inputs}, not over a fixed final circuit. Given a program subject---a builder plus the auditable semantics of its inputs---QMTester derives a follow-up \emph{input}, rebuilds the circuit from that follow-up input, executes the source and follow-up circuits under matched settings, canonicalizes the measured count vectors through a validated bijective map, and applies a corrected two-sample test to decide whether the program-level relation is violated. The relation families capture builder-level contracts: qubit-role covariance under input permutation, classical-register remapping, QFT round-trip direction, parameter periodicity, and ancilla uncompute. Each relation is admitted only under explicit, documented preconditions; a candidate that would change the number of measured bits, drop a register, or induce a non-bijective output map is rejected rather than tested.

We are deliberate about scope: we do not claim benchmark dominance over Bugs4Q, and the injected control is not evidence of real-world generality. Our calibrated, auditable claim is that on defects with a documented sound relation, program-level testing \emph{matches a ground-truth reference-oracle differential with no oracle}---not that it beats MorphQ, whose $0/7$ is structural blindness to the builder layer, not a comparative loss---at zero false positives on the matched fixed variants. A circuit-level-only diagnostic in our own artifact shows \emph{no} advantage over MorphQ, motivating the move to the program level rather than a larger circuit-level catalog.

The empirical account is traceable: every headline number resolves to a row of a canonical CSV derived from frozen manifests and run-isolated shards, and each included subject's relation-soundness audit is documented for inspection.

We make three contributions:
\begin{itemize}
  \item \textbf{A program-level instantiation of metamorphic testing for Qiskit.} We argue and demonstrate that a quantum program is a builder, parameterized APIs, register mappings, and ancilla protocols---not only its final circuit---and that defects in these contracts are a distinct target. For a \emph{builder-ignores-input} defect, no circuit-level metamorphic self-transform (the MorphQ class) can produce a source/follow-up divergence---it rewrites the already-emitted buggy circuit rather than re-invoking the builder on a different input---so detection requires rebuilding from a follow-up input (Proposition~1); this necessity holds structurally for qubit-role permutation and classical-register remapping, and across the non-unitary boundary (\S\ref{sec:nonunitary}) for ancilla-uncompute, whereas \texttt{program\_qft\_round\_trip} and \texttt{program\_parameter\_periodicity} are circuit-expressible and included for coverage.
  \item \textbf{Canonicalization-as-admission, in five distribution families plus a complementary static check.} The technical core is an admission discipline: a relation is tested only when its output map is a \emph{documented bijection on the measured support}, which separates a representation change from a logical divergence---candidates that drop a register or yield a non-bijective map are rejected, not scored. We instantiate it in five \texttt{program\_*} families (\texttt{input\_permutation}, \texttt{classical\_remap}, \texttt{qft\_round\_trip}, \texttt{parameter\_periodicity}, \texttt{ancilla\_uncompute}) under a quantum-aware period check (a $2\pi$ shift is rejected on a controlled rotation, whose period is $4\pi$), with a \emph{complementary static} \texttt{program\_register\_shape} check---needing no statistics---that brings the prevalent \texttt{measure\_all} register-shape defect in scope.
  \item \textbf{A prevalence finding (first-class), and an inclusion-independent evaluation.} An exhaustive screen of two real-defect corpora characterizes the niche precisely: distribution-detectable builder defects are \emph{rare}---the most common builder defect is the \texttt{measure\_all} register-shape change ($37$ of LintQ's $33{,}995$ findings), caught by the static register-shape check---so with all 32 build-and-run Bugs4Q bugs audited (9 in scope), we report a real-but-small niche rather than overstate it. Detection on the in-scope defects is \emph{deterministic} for all four real families (TVD~$=1$, the real QFT subject included) and holds $7/7$ under a realistic readout/crosstalk noise model; the statistics are load-bearing for specificity and the one probabilistic family, not for detection. The inclusion-\emph{independent} evidence is specificity (0/107 fixed variants, 0/50 correct programs; smallest per-pair $p=0.0035$ under the realistic model, Bonferroni still 0/50) and a detectability-unfiltered power curve (floor near TVD~$0.025$), plus an oracle-light ablation.
\end{itemize}
